\section{Week 13}

\subsection{Automorphisms}

Let \( G  \) be a group and consider \( {S}_{G} \), not all of the permutations in \( {S}_{G} \) preserve algebraic information about \( G \). We consider the class of permutations of \( G  \) which are isomorphisms. Consider an isomorphism. Let's call the mapping \( \varphi : G \to G  \). Since \( \varphi  \) is a bijection, \( \varphi  \) is a permutation of \( G  \). That is, \( \varphi \in {S}_{G} \). Since \( \varphi  \) is an isomorphism, it also preserves the operation of \( G  \) when it maps by definition.  

\begin{definition}[ ]
    Let \( G  \) be a group. An isomorphism from \( G  \) onto itself is called an \textbf{automorphism of \( G  \)}. The set of all automorphisms of \( G  \) is denoted as \( \text{Aut}(G) \).
\end{definition}

\begin{remark}
   \begin{enumerate}
       \item[(1)] \( \text{Aut}(G) \) form a group with respect to composition. 
        \item[(2)] \( \text{Aut}(G) \leq {S}_{G} \) 
        \item[(3)] A special kind of automorphism is \textbf{conjugation}.
   \end{enumerate} 
\end{remark}

\begin{prop}[Proposition 13 from the book]
    Let \( H \trianglelefteq G  \). 
    \begin{enumerate}
        \item[(1)] \( G  \) acts on \( H  \) by conjugation (\( g \cdot h = g h g^{-1} \) for all \( g \in G  \) and \( h \in H  \) is a group action);
        \item[(2)] For \( g \in G  \), the mapping \( {\varphi}_{g} : H \to H  \) such that \( {\varphi}_{g}(h) = g h g^{-1} \) for all \( h \in H  \) is an automorphism of \( H  \);
        \item[(3)] \( G / {C}_{G}(H) \cong K \leq \text{Aut}(H) \).
    \end{enumerate}
\end{prop}

\begin{proof}
\begin{enumerate}
    \item[(1)] Defining the action as in the statement, we note that the two axioms of a group action follow immediately from the operation in \( G  \) (specifically the identity of \( G  \) and associativity. To see that this defines a map from \( G \times H  \) to \( H  \). Note that \( H \trianglelefteq G  \) and so \( g h g^{-1} \in H   \) for all \( g \in G  \) and \( h \in H  \).   
    \item[(2)] We need to show that \( {\varphi}_{g} : H \to H  \) defined by \( {\varphi}_{g}(h) = g h g^{-1}\) is a bijection. Since \( H \trianglelefteq G  \), \( g H g^{-1} = H  \) and hence, \( h \mapsto g h g^{-1} \) is a bijection from \( H  \) to \( H  \). (3.2 Problem 5 from Dummit and Foote)
Note that \( {h}_{1}, {h}_{2} \in H  \) by closure and so \( {\varphi}_{g}  \) is an automorphism of \( H  \). Further for \( {h}_{1}, {h}_{2} \in H \),
        \[  {\varphi}_{g}({h}_{1}) {\varphi}_{g}({h}_{2})  = g {h}_{1} (g^{-1} g ) {h}_{2} g^{-1}  =  g {h}_{1} {h}_{2} g^{-1} = {\varphi}_{g}({h}_{1}{h}_{2}). \]

    \item[(3)]      Letting \( \psi  \) be the permutation representation of \( G  \) acting on \( H  \) via conjugation, we see that 
    \[  \psi: G \to {S}_{H} \]
    defined by \( \psi(g) = {\varphi}_{g} \)
    where \( {\varphi}_{g}(h) = g h g^{-1} \). By the first isomorphism theorem, we have  
    \[  G / \text{ker}(\psi) \cong \psi(G) \leq {S}_{H}. \]
    Since \( \psi(g) \) is an automorphism of \( H  \) for all \( g \in G  \), we have that 
    \[  \psi(G) \leq \text{Aut}(G). \]
    Lastly, the kernel of \( \psi  \) is equal to
    \begin{align*}
        \text{ker}(\varphi) &= \{  g \in G : \psi(g) = {1}_{{S}_{H}} \}  \\
                            &= \{ g \in G : {\varphi}_{g}(h) = h \ \forall h \in H  \} \\ 
                            &= \{  g \in G : g h g^{-1} = h \ \forall h \in H  \} \\
                            &= {C}_{G}(H).
    \end{align*}
\end{enumerate}
\end{proof}

\begin{corollary}
    If \( K \leq G  \), then \( K \cong g K g^{-1} \) and conjugate elements and conjugate subgroups have the same order.  
\end{corollary}
\begin{proof}
If we let \( H = G  \), then \( {\varphi}_{g} : G \to G  \) is an isomorphism so the result follows. 
\end{proof}

\begin{corollary}
    For any \( H \leq G  \), we have 
    \(  {N}_{G}(H) / {C}_{G}(H)  \)
    is isomorphic ot a subgroup of \( \text{Aut}(H) \). In particular, \( G / Z(G) \cong K \leq \text{Aut}(G)  \).
\end{corollary}
\begin{proof}
We know that \( H \trianglelefteq {N}_{G}(H) \). Letting \( G  \) be \( {N}_{G}(H) \) in our first theorem gives us the first statement. For the second statement, when \( H = G  \) in our initial theorem. We have \( G / Z(G) \cong K \leq \text{Aut}(G) \) since \( {C}_{G}(G) = Z(G) \).
\end{proof}

\begin{definition}[Inner Automorphisms]
    Let \( G  \) be a group and \( g \in G  \). Conjugation by \( g  \) (\( {\varphi}_{g} \)) is called an \textbf{inner automorphism of \( G  \)}.
The subgroup of \( \text{Aut}(G) \) consisting of all inner automorphisms of \( G \) is denoted \( \text{Inn}(G) \).
\end{definition}

Note that \( \text{Inn}(G) = \{  {\varphi}_{g} : g \in G  \}  \) where \( {\varphi}_{g}(x) = g x g^{-1}\) for all \( x \in G  \). By our 2nd corollary, we have that 
\[  G / Z(G) \cong \psi(G) = \{ {\varphi}_{g} : g \in G \} = \text{Inn}(G). \]

\begin{eg}
    \begin{enumerate}
        \item[(1)] A group is abelian if and only if every inner automorphism is trivial (if \( g x g^{-1} = x  \) for all \( x,g \in G  \))
        \item[(2)] Let \( G = {D}_{8} \) and \( Z({D}_{8} ) = \{  1, r^{2} \}  = \langle  r^{2} \rangle \) where 
            \[  {D}_{8} / Z({D}_{8}) = {D}_{8} / \langle  r^{2} \rangle \cong \text{Inn}({D}_{8}). \]
            If we remember our lattice tree for \( {D}_{8} \). Since \( | {D}_{8} / \langle  r^{2} \rangle |  =  4  \), we know that  
            \[  {D}_{8}/ \langle  r^{2} \rangle \cong {V}_{4} \ \text{or} \ {Z}_{4}. \]
            But we know that \( {D}_{8} / \langle  r^{2} \rangle \) is not cyclic and so it must be isomorphic to \( {V}_{4}  \) (the Klein \( 4  \) group).
        \item Let \( G = {Q}_{8}  \) and \( Z({Q}_{8}) = \{  \pm 1  \}  \) and we can perform a similar process to the previous example with the same structure as the Klein group.
        \item Since \( Z({S}_{n}) \) where \( n \geq 3  \), the class equation says that 
        \[  | {S}_{n} |  = | Z(G) |  + \sum (\text{conjugacy class not in the center}). \]
    \end{enumerate}
    Recall that the conjugates of \( {S}_{n} \) are just its cycle types.
    The only element of \( {S}_{n} \) which has no other conjugates is the identity and so \( Z({S}_{n}) = {1}_{{S}_{n}} \). Then from our theorem 
    \[  {S}_{n} / Z({S}_{n}) \cong \text{Aut}({S}_{n}). \]
\end{eg}

\subsection{What does \( \text{Aut}(G) \)} 

\begin{prop}
    The automorphism group of a cyclic group of order \( n  \) has \( \varphi(n) \) where (\( \varphi  \) is the Euler-Totient function).
\end{prop}

\begin{proof}
    Suppose \( G = \langle x  \rangle \) and \( | x  |  = n < \infty   \) and \( \varphi  \) is an automorphism of \( G  \). Then \( \varphi(x) = x^{a}  \) for some \( 1 \leq a \leq n  \). We also know 
    \[  | x^{a} |  = \frac{ n  }{  \text{gcd}(n,a)  }  \ \ \text{so} \ \ | x^{a} |  = n  \iff \text{gcd}(a,n) = 1.   \]
    Hence, \( \varphi  \) can send \( x  \) to a positive integer less than or equal to \( n  \) which is relatively prime. Hence, there are \( \varphi(n) \) choices. 
\end{proof}

\begin{eg}
    If \( G  \) is cyclic of order \( 5^{2} \cdot 2^{3} \) then 
    \begin{align*}
      | \text{Aut}(G)  |  &= \varphi(5^{2} \cdot 2^{3}) \\
                          &= \varphi(5^{2})\varphi(2^{3}) \\
                          &= 5(5 - 1) \cdot 2^{2}(2 - 1) \\
                          &= 5 \cdot 4 \cdot 4 \\
                          &= 80.
    \end{align*}
\end{eg}

\begin{prop}
    Assume \(G  \) is a group of order \( p q  \) where \( p  \) and \( q  \) (not necessarily distinct) with \( p < q  \) if \( p \not\mid q - 1  \), then \( G  \) is abelian. 
\end{prop}

\begin{proof}
We will consider two cases
\end{proof}

